In rocksDB, range deletes hold information about non-existent entries, but they are solely stored in files, resulting in additional disk IO operations in accessing them, inevitably declining query speed. Presently, the bloom filter only stores details about point entries, and fencing pointers store information about existing file ranges. However, there is no filter that stores delete ranges, making it impossible to determine whether a specific point is deleted without resorting to disk IO. \\

First, current system loss a degree of freedom for not keeping range delete information in the buffer. One advantage of rocksDB is that users can choose to keep the bloom filter inside the buffer, enabling them to filter out non-existent keys in a given file without needing direct disk access. Unfortunately, this option is not available for range deletes. When a user wants to determine if a key is deleted by a range delete or not, performing disk IO operations becomes necessary. Consequently, we lose the freedom to remove range-deleted keys before directly accessing the disk, leading to potential impacts on performance. \\

In Addition, without range delte filter, systems take more time in reading an SST file. For a system, files are read in blocks and each file contains multiple blocks. To read an SST file, rocksDB first reads the metadata to identify which block may contain the target key or have overlapping deleted ranges with it. Then, it proceeds to find the key in the block that may contain the key and retrieve range delete information in range delete block. Since range deletes are currently not stored in the buffer, the system must perform extra disk IOs to navigate through the range delete blocks to determine whether a key is range-deleted. \\